% Figure 11.2 : deux conteneurs d'un même Pod, vus par le runtime (inodes réels du chapitre 11).
\documentclass[tikz,border=4pt]{standalone}
\usepackage{quai}
\begin{document}
\begin{tikzpicture}[x=1cm,y=1cm,
    ns/.style={qbox, font=\ttfamily\scriptsize, minimum width=2.6cm, minimum height=0.62cm, inner sep=2pt},
    part/.style={ns, draw=QSarcelle, fill=QSarcelleBg},
    propre/.style={ns, draw=QGris, fill=QGrisBg}]

  \node[qcadre, minimum width=10.8cm, minimum height=5.9cm, draw=QOcre] at (2.9,1.15) {};
  \node[font=\titrefig\normalsize, text=QOcre, anchor=north west] at (-2.4,4.0) {Pod demo-runtime};
  \node[font=\ttfamily\scriptsize, text=QMuted, anchor=north east] at (8.2,4.0) {containerd-shim, PID 4022};

  \node[font=\titrefig\normalsize] at (0,3.05) {pause (PID 4047)};
  \node[font=\titrefig\normalsize] at (5.8,3.05) {nginx (PID 4071)};

  \foreach \y/\t/\i in {2.2/net/4026534223, 1.35/uts/4026534287, 0.5/ipc/4026534288} {
    \node[part] (p\t) at (0,\y) {\t:[\i]};
    \node[part] (n\t) at (5.8,\y) {\t:[\i]};
    \draw[qarr, draw=QSarcelle] (n\t.west) -- (p\t.east);
  }
  \node[propre] at (0,-0.35) {pid:[4026534289]};
  \node[propre] at (5.8,-0.35) {pid:[4026534291]};
  \node[propre] at (0,-1.2) {mnt:[4026534286]};
  \node[propre] at (5.8,-1.2) {mnt:[4026534290]};

  \node[qlblc=QSarcelle, font=\scriptsize, fill=QPaper, inner sep=1.5pt] at (2.9,2.62) {rejoint, par {\ttfamily /proc/4047/ns/…}};
  \node[qlbl, font=\scriptsize, fill=QPaper, inner sep=1.5pt] at (2.9,-0.77) {créés pour chacun};
\end{tikzpicture}
\end{document}
